\documentclass[../main.tex]{subfiles} % 使用 subfiles 文档类，指定主文档
\graphicspath{{\subfix{../image/}}} % 指定图片目录，后续可以直接使用图片文件名。

% 例如:
% \begin{figure}[H]
% \centering
% \includegraphics[width=0.3\textwidth]{图.png}
% \caption{}
% \label{figure:图}
% \end{figure}
% 注意:上述\label{}一定要放在\caption{}之后,否则引用图片序号会只会显示??.

\begin{document}
% SOURCE p=128; printed=113; object=ch04-title; region=top
\chapter{点集拓扑}\label{chapter:ch04-topology}
% SOURCE p=128; printed=113; object=ch04-introduction; region=middle
极限、收敛和连续性是整个分析学的核心概念。为它们建立一个把经典情形包含为特例的一般框架，是很有用的。度量空间提供了这样一个框架，其优点是：在 Euclid 空间分析中已有的概念之外，不需要引入很多新概念。然而，度量空间还不够一般，甚至不能描述某些十分经典的收敛方式，例如 $\mathbb R$ 上函数的逐点收敛。把开集而不是度量作为原始数据，可以建立一种更灵活的理论；本章将研究的正是这一理论。
\subfile{Chapter-04/section-01.tex}
\subfile{Chapter-04/section-02.tex}
\subfile{Chapter-04/section-03.tex}
\subfile{Chapter-04/section-04.tex}
\subfile{Chapter-04/section-05.tex}
\subfile{Chapter-04/section-06.tex}
\subfile{Chapter-04/section-07.tex}
\subfile{Chapter-04/section-08.tex}
\subfile{Chapter-04/section-09.tex}
\end{document}
